\notepage{15a}{Simultaneous diagonalization of two Hermitian operators}
\smallheading{Theorem 13}
If $\Omega$ and $\Lambda$ are Hermitian and commute,
\[
 \comm{\Omega}{\Lambda}=0,
\]
there exists a complete orthonormal basis of simultaneous eigenvectors:
\[
 \Omega\ket{\omega_i}=\omega_i\ket{\omega_i},\qquad
 \Lambda\ket{\omega_i}=\lambda_i\ket{\omega_i}.
\]

\smallheading{Why each eigenspace is preserved}
Starting with $\Omega\ket{\omega_i}=\omega_i\ket{\omega_i}$, multiply
by $\Lambda$ and commute the two operators:
\begin{align*}
 \Lambda\Omega\ket{\omega_i}&=\omega_i\Lambda\ket{\omega_i},\\
 \Omega\bigl(\Lambda\ket{\omega_i}\bigr)
 &=\omega_i\bigl(\Lambda\ket{\omega_i}\bigr).
\end{align*}
Thus $\Lambda\ket{\omega_i}$ remains in the $\omega_i$ eigenspace
(it may be zero).
\begin{enumerate}[label=(\roman*)]
\item If $\omega_i$ is nondegenerate, that eigenspace is
one-dimensional, so $\Lambda\ket{\omega_i}=\lambda_i\ket{\omega_i}$.
\item If $\omega_i$ is degenerate, $\Lambda$ can mix vectors inside
that eigenspace, but cannot mix it with an eigenspace of different
$\omega$. In particular,
\[
 \omega_i\ne\omega_j\quad\Longrightarrow\quad
 \mel{\omega_j}{\Lambda}{\omega_i}=0.
\]
\end{enumerate}

\smallheading{Block structure in an eigenbasis of $\Omega$}
Group the equal eigenvalues together. The handwritten block sketches
can be written as
\[
 [\Omega]=\begin{bmatrix}
 \omega_1 I_{m_1}&0&\cdots\\
 0&\omega_2 I_{m_2}&\cdots\\
 \vdots&\vdots&\ddots
 \end{bmatrix},\qquad
 [\Lambda]=\begin{bmatrix}
 \Lambda_1&0&\cdots\\
 0&\Lambda_2&\cdots\\
 \vdots&\vdots&\ddots
 \end{bmatrix}.
\]
Each $\Lambda_a$ acts only inside the corresponding
$\omega_a$ eigenspace. The off-block matrix elements vanish.
